\subsection{Prospective scenario validation}
\label{subsec:orf-prospective}
After freezing the 100\,m\textsuperscript{3} policy, we generated 24 new scenarios
at 24, 72 and 168 hourly periods, using fixed random windows and 10\% perturbations
of inflow and demand from the six known clean inputs. These are new parameter
scenarios within one family, not independent operational data. Five policies and
three new seeds give 360 interleaved calls with the previous solver settings.
Table~\ref{tab:orf-prospective-main} brings the precision and effort outcomes
together. All five policies meet the budgets in 72/72 assignments, so the
contract provides no additional primal feasibility recovery in this sample.
Exact witnesses detect seven L2 bound contradictions, one Budget-L2 contradiction
and none for L2-NF3. Satisfied primal budgets do not certify solver bounds.

\begin{table}[htbp]
\centering\small\setlength{\tabcolsep}{3pt}
\caption{Prospective validation on 24 scenarios with three seeds each. Budgets
count passing runs; B/G count detected bound/1\%-completion contradictions.
Work and complete-process time are paired geometric means for Budget-L2 divided
by the row policy, with marginal 95\% scenario-cluster bootstrap intervals below.
Wins are scenarios favoring Budget-L2 after aggregating their three seeds,
for work/time respectively, out of 24. All adverse outcomes are included.}
\label{tab:orf-prospective-main}
\begin{tabular}{lccccc}\toprule
Policy & Budgets & B/G & Work ratio & Time ratio & Wins $W/T$\\\midrule
Base & 72/72 & 0/0 & \shortstack{0.422\\{[0.288, 0.604]}} & \shortstack{0.797\\{[0.646, 0.961]}} & 21/16\\[4pt]
Base-tight & 72/72 & 0/0 & \shortstack{0.170\\{[0.136, 0.213]}} & \shortstack{0.643\\{[0.524, 0.772]}} & 24/21\\[4pt]
L2 & 72/72 & 7/6 & \shortstack{0.931\\{[0.702, 1.133]}} & \shortstack{0.996\\{[0.832, 1.119]}} & 10/6\\[4pt]
L2-NF3 & 72/72 & 0/0 & \shortstack{0.476\\{[0.377, 0.582]}} & \shortstack{1.031\\{[0.935, 1.129]}} & 24/8\\[4pt]
Budget-L2 & 72/72 & 1/0 & --- & --- & ---\\[4pt]
\bottomrule\end{tabular}\end{table}

The uncertainty analysis was added after the experimental outcomes were known;
it was not part of the frozen optimization protocol. Following the cluster
bootstrap principle of retaining within-group observations~\citep{FieldWelsh2007},
we resample 24 scenarios with replacement 100,000 times, keeping all three seeds
and paired policies together. We average log ratios within each scenario and
then across scenarios; the balanced design gives the same point estimate as
the 72-pair geometric mean. Percentile intervals are conditional on this
generated sample and its six known parents. They do not assess generalization
to other applications or account for selecting different parent structures.

Against L2, Budget-L2 lowers work in 24/72 seed pairs and 10/24 scenario
aggregates; it lowers complete-process time in 6/24 scenarios. The work ratio
is 0.931 with interval [0.702,1.133], and the time ratio is 0.996 with interval
[0.832,1.119]. Both intervals include one. Horizon-specific work ratios are
0.944, 1.053 and 0.811 for 24, 72 and 168 hours. Eight scenarios per horizon
cannot establish a structural selection rule. A geometric mean below one
can coexist with a minority of wins because improvement magnitudes vary.
After excluding pairs with contradictions in either policy, the L2 comparison
is 0.971 for work and 1.023 for time (64 pairs); this secondary screen does not
certify the retained bounds.

The incremental comparison is narrower than against Base: L2 already changes
the representation using the same Euclidean kernel, whereas the contract
modifies only factors required by the precision specification (8,448 activation
rows across these models). This identifies the intervention without proving
its internal solver mechanism. The evidence supports explicit specification,
admissible selection and verification of accuracy, with conditional numerical
benefits; it does not establish general efficiency gains over L2.

Budget-L2's maximum candidate cost deterioration against Base is 1.020\%,
exceeding the earlier six-model maximum below 0.1\%. This is the paired quantity
$100(C_{\rm Budget}-C_{\rm Base})/|C_{\rm Base}|$, not a difference between
certified MILP optima. The requested 1\% solver gap neither bounds this
cross-policy candidate difference nor establishes equally certified gaps,
especially when an independent witness contradicts a reported bound.
Online Resource~1 reports every stratum, intervention and rational-reference
outcome, together with the registered 360-call logging confirmation.
